\newpage
\section{Generalidades del proyecto}
%----------------------------------------------------------------------------

\subsection{Descripción del problema}
La reventa de entradas en canales digitales enfrenta problemas de autenticidad, trazabilidad y transparencia de transacciones. En plataformas no especializadas y en mercados informales, el comprador no siempre cuenta con mecanismos verificables para confirmar propiedad, historial del ticket y validez al momento de ingreso.


\subsection{Objetivos}
\subsubsection{Objetivo General}
Desarrollar una solución B2B de reventa de tickets, basada en contratos inteligentes de Stellar, que garantice unicidad y autenticidad del activo digital y permita evaluar mecanismos de reducción de fraude y especulación en el mercado secundario.
%----------------------------------------------------------------------------

\subsubsection{Objetivos Específicos}
 
\begin{itemize}
    \item Diseñar un modelo de ticket NFT con unicidad verificable y trazabilidad de propiedad.
    \item Implementar contratos inteligentes para emisión, listado, compra, transferencia y redención de tickets.
    \item Construir un prototipo web distribuido para simular el flujo completo de reventa en testnet.
    \item Evaluar desempeño técnico y comportamiento antifraude mediante métricas reproducibles.
\end{itemize}
%----------------------------------------------------------------------------

\subsection{Alcance y delimitaciones}

El alcance se centra en el mercado secundario de entradas con despliegue en testnet. Quedan fuera de alcance, para esta versión, integración con pasarelas fiat reales, KYC productivo, despliegue definitivo en mainnet y acoplamiento con tiqueteras comerciales de terceros.

\subsection{Metodología (investigación + desarrollo del prototipo)}

La metodología combina revisión técnico-científica, diseño de arquitectura, implementación incremental del prototipo y evaluación experimental. Se articula en ciclos de definición de requerimientos, desarrollo de componentes y validación por escenarios de uso y fraude.

\subsection{Requerimientos}

\subsubsection{Requerimientos funcionales (mínimos: 10)}

Incluyen, como base, creación y gestión de tickets, listado de reventa, compra atómica, distribución de comisiones, verificación de titularidad, redención en puerta y trazabilidad de transferencias.

\subsubsection{Requerimientos no funcionales (mínimos: 6, base ISO/IEC 27001)}

Consideran control de acceso, integridad de datos, trazabilidad de eventos, gestión de claves, disponibilidad operativa y monitoreo de fallos en componentes críticos.

\subsubsection{Requerimientos complementarios (no obligatorios)}

Como extensiones se consideran reglas dinámicas de mercado, políticas avanzadas de cancelación, analítica de especulación y soporte para activos de pago adicionales.

\subsection{Estado del arte (comparación técnica y evidencia empírica)}

La literatura y la evidencia de mercado reportan riesgos persistentes de fraude en ticketing digital, especialmente en contextos de alta demanda. Como línea comparativa, este trabajo contrasta mecanismos tradicionales de reventa con enfoques basados en blockchain en términos de trazabilidad, verificabilidad y tiempos de liquidación.

\subsection{Marco legal colombiano (caso de referencia: Ley 1493/2011 y fiscalización digital)}

El proyecto usa la Ley 1493 de 2011 como marco de referencia documental para discutir fiscalización y trazabilidad digital de transacciones en espectáculos públicos. El análisis jurídico se modela como mapeo normativo simulado dentro del diseño, sin declarar cumplimiento regulatorio productivo.
\subsubsection{Requerimientos de conectividad}

